\begin{satz}{Fortsetzung einer endlichen Erwartungswertungleichung (I)}
            {AbzaehlbareDarstellungUndInfimumVonSuprema}
Seien $I, J$ nichtleere, abz�hlbare Mengen,
$X, Y: I\times J \rightarrow \mcm(\Omega, \mfa, \R)$, sowie
$\mci := \Pot_{ne}(I)$ und $\mcj := \Pot_{ne}(J)$.\\
Es gelte:
\begin{enumerate}[(a)]
\item $\forall\,\omega\in\Omega\,\forall\, Z \in \meng{X, Y}: 
      \menge{Z(i, j)(\omega)}{(i, j)\in I\times J}$
      beschr�nkt.
\item $\forall\, A \in \mci \cup \meng I \, \forall\, B \in \mcj \cup \meng J:
     \Inf{i \in A}\Sup{j \in B} X(i, j) = \Sup{j \in B}\Inf{i \in A} X(i,j)$.
\item $\forall A \in \mci \cup \meng I  \forall B \in \mcj \cup \meng J:
     \Inf{i \in A}\Sup{j \in B} X(i, j), \Inf{i \in A}\Sup{j \in B} Y(i, j)
     \in L_1(\Omega, \mfa, \Prim, \R)$.   
\item $\forall\, A \in \mci \, \forall\, B \in \mcj:
\E\bigg(\underset{i \in A}\inf\,\underset{j \in B}\sup\, X(i, j) \bigg)
\le \E\bigg(\underset{i \in A}\inf\,\underset{j \in B}\sup\, Y(i, j) \bigg)$.\\
(Alle genannten Abbildungen sind wegen $(a)$ und $I \neq \emptyset \neq J$
 wohldefiniert.)
\end{enumerate}
Dann gilt \[\E\left(\inf_{i \in I}\,\sup_{j \in J}\, X(i, j) \right)
\le \E\left(\inf_{i \in I}\,\sup_{j \in J}\, Y(i, j) \right)\text.\]
\end{satz}
\begin{proof}
Definiere $W_j := \Inf{i\in I} X(i, j)$ f�r alle $j \in J$.\\
Wegen (a) ist $\menge{W_j(\omega)}{j\in J}$ f�r alle $\omega \in \Omega$ nach
oben beschr�nkt.\\
F�r jedes $k\in J$ ist $W_k$ und
\begin{align}
\sup_{j \in J}W_j = \sup_{j \in J}\inf_{i \in I}X(i, j) \overset{(b)}= 
\inf_{i \in I}\sup_{j \in J}X(i, j)\quad \text{ nach (c) integrierbar.}
\label{feeeieq:1}
\end{align}
Damit sind die Voraussetzungen von
\ref{EndlicheDarstellungDesSupremumsIntegrierbarerFunktionen} f�r
$\folge W j J$ erf�llt. Daher ist
\[\menge{\E\left(\inf_{i\in I}\sup_{j \in B}X(i,j) \right)}{B \in \mcj}
\overset{\eqref{feeeieq:1}}= 
 \menge{\E\left(\sup_{j \in B}W_j \right)}{B \in \mcj}\]
nach \refclaim{EndlicheDarstellungDesSupremumsIntegrierbarerFunktionen}
{EndlicheDarstellungDesSupremumsIntegrierbarerFunktionen1} nichtleer und nach
oben beschr�nkt und es gilt:
\begin{align}
& \sup\menge{\E\left(\inf_{i \in I}\sup_{j \in B}X(i, j) \right)}{B \in \mcj}
\notag\\
 =&\,  \sup\menge{\E\left(\sup_{j \in B}W_j \right)}{B \in \mcj}
 \overset{\refclaim{EndlicheDarstellungDesSupremumsIntegrierbarerFunktionen}
 {EndlicheDarstellungDesSupremumsIntegrierbarerFunktionen2}}=
\E\left(\sup_{j \in J}W_j \right)\notag\\ 
\overset{(b)}= &\,\E\left(\inf_{i\in I}\sup_{j\in J}X(i,j)\right)\text.
\label{feeeieq:2}
\end{align}
Seien nun $Z_i := \Sup{j \in J}Y(i, j)$ f�r alle $i \in I$.\\
Wegen (a) ist $\menge{Z_i(\omega)}{i\in I}$ nach unten beschr�nkt f�r alle
$\omega \in \Omega$.\\
F�r alle $l\in I$ ist $Z_l$ und
\[\inf_{i \in I}Z_i = \inf_{i \in I}\sup_{j \in J}Y(i, j) \quad
\text{ nach (c) integrierbar.} \]
Damit sind auch die Vorausssetzungen von
\ref{EndlicheDarstellungDesInfimumsIntegrierbarerFunktionen} erf�llt. Also ist
\[\menge{\E\left(\inf_{i\in A}\sup_{j\in J}Y(i,j)\right)}{A\in\mci}
  =\menge{\E\left(\inf_{i\in A}Z_i\right)}{A\in\mci}\]
nach \refclaim{EndlicheDarstellungDesInfimumsIntegrierbarerFunktionen}
{EndlicheDarstellungDesInfimumsIntegrierbarerFunktionen1} nichtleer und nach
oben beschr�nkt und es gilt:
\begin{align}
& \inf\menge{\E\left(\inf_{i\in A}\sup_{j\in J}Y(i,j)\right)}{A\in\mci}
= \inf\menge{\E\left(\inf_{i \in A}Z_i \right)}{A \in \mci}\notag\\
\overset{\refclaim{EndlicheDarstellungDesInfimumsIntegrierbarerFunktionen}
{EndlicheDarstellungDesInfimumsIntegrierbarerFunktionen2}}=&
\E\left(\inf_{i\in I}Z_i\right)=\E\left(\inf_{i\in I}\sup_{j\in J}Y(i,j)\right)
\text.\label{feeeieq:3}
\end{align}
Seien nun $C \in \mci$ und $D \in \mcj$. Dann gilt
\[ \inf_{i \in I}\sup_{j \in D}X(i, j) \le \inf_{i \in C}\sup_{j \in D}X(i,j)
\quad \text{ und }\quad  
\inf_{i\in C}\sup_{j\in D}Y(i,j)\le \inf_{i \in C}\sup_{j \in J}Y(i, j)\]
und diese vier Abbildungen sind nach (c) integrierbar.\\
Daher gelten diese Ungleichungen auch im Erwartungswert, sodass gilt:
\begin{align*}
\E\left(\inf_{i \in I}\sup_{j \in D}X(i, j) \right) & \le
\E\left(\inf_{i \in C}\sup_{j \in D}X(i, j) \right) \overset{(d)}\le
\E\left(\inf_{i \in C}\sup_{j \in D}Y(i, j) \right)\\
&\le \E\left(\inf_{i \in C}\sup_{j \in J}Y(i, j) \right)\text.
\end{align*}
Somit erh�lt man
\begin{align}
&
\sup\menge{\E\left(\inf_{i\in I}\sup_{j\in B}X(i,j)\right)}{B \in \mcj}\notag\\
\le&\, \inf\menge{\E\left(\inf_{i\in A}\sup_{j\in J}Y(i,j)\right)}{A\in\mci}
 \text. \label{feeeieq:4}\end{align}
Dies liefert schlie�lich
\begin{align*}
& \E\left(\inf_{i\in I}\sup_{j\in J}X(i,j)\right) \overset{\eqref{feeeieq:2}}=
\sup\menge{\E\left(\inf_{i \in I}\sup_{j \in B}X(i, j) \right)}{B \in \mcj}\\
\overset{\eqref{feeeieq:4}}\le &
\inf\menge{\E\left(\inf_{i\in A}\sup_{j\in J}Y(i,j)\right)}{A\in\mci}
\overset{\eqref{feeeieq:3}}= \E\left(\inf_{i\in I}\sup_{j\in J}Y(i,j)\right)
\end{align*}
und damit die Behauptung.
\end{proof}